\documentclass[10pt,a4paper]{amsart}
\usepackage[margin=19mm]{geometry}
\usepackage{amsmath,amssymb,mathtools}
\usepackage[scr=rsfs,cal=euler]{mathalfa}
\usepackage{xfrac}
\usepackage[unicode,hidelinks,bookmarks=false]{hyperref}
\newcommand{\ie}{i.e.\ }
\newcommand{\cf}{cf.\ }
\newcommand{\resp}{respectively\ }
\newcommand{\Subsection}{\subsection}
\providecommand{\nobreakdash}{\nobreak\hbox{-}}
\setlength{\emergencystretch}{2em}
\multlinegap0pt
\usepackage{bm}
\usepackage{bbm}
\usepackage{dsfont}
\usepackage{xspace}
\usepackage{cancel}
\usepackage{mathtools}
\usepackage[shortlabels]{enumitem}
\usepackage{amsthm}
\makeatletter
\@ifundefined{theorem}{\newtheorem{theorem}{Theorem}}{}
\@ifundefined{proposition}{\newtheorem{proposition}[theorem]{Proposition}}{}
\makeatother
\newcommand{\aut}{\mathfrak{aut}}
\newcommand{\calC}{\mathcal{C}}
\newcommand{\calL}{\mathcal{L}}
\newcommand{\hor}{\mathrm{H}}
\newcommand{\pd}[2]{\frac{\partial #1}{\partial #2}}
\newcommand{\ver}{\mathrm{V}}
\renewcommand{\d}{\mathrm{d}}
\title[Symmetries of one-loop deformed q-map spaces]{Symmetries of one-loop deformed q-map spaces}
\author{Vicente Cort\'es, Alejandro Gil-Garc\'ia, Danu Thung}
\date{Source: arXiv:2402.16178v1 (CC BY 4.0)}
\begin{document}
\maketitle

\noindent\emph{Source and licence.} Symmetries of one-loop deformed q-map spaces. Version arXiv:2402.16178v1 (CC BY 4.0), distributed under the Creative Commons Attribution 4.0 International licence, as recorded in the arXiv licence metadata for this exact version and shown on the abstract page https://arxiv.org/abs/2402.16178v1. Full rights notice: licenses/arxiv-2402.16178-CC-BY-4.0.txt.\par\smallskip

\noindent\textit{Editorial scope.} Counted unit: the Proposition whose source label is \texttt{prop:moment\_map} in the article (source0.tex lines 421--423, source label \texttt{prop:moment\_map}), delivered together with the article's own complete original proof of it (source0.tex lines 425--448, one \texttt{proof} environment of 24 lines). No mathematics is written, repaired or re-derived: statement and proof are the article's own text line for line. Required context retained uncounted: prop:basis\_CASK (lines 311--318); oH:eq (lines 388--393); pi*nabla:eq (lines 413--415). Every definition, display and label of those lines is retained unchanged. Omitted from this package and not redistributed: 1--310, 319--387, 394--412, 416--420, 424--424, 449--645. Nothing inside a retained excerpt is omitted or altered: every retained body is delivered exactly as the article's lines are written, so no omission receipt of any kind accompanies this package. Citations: the retained excerpts carry no citation command at all, so no bibliography entry of the article is transcribed and no reference is redistributed. Reference adaptation: none was needed and none was made; every reference inside the delivered unit resolves inside it, so no unresolved reference or citation is produced. Printed slips kept literal: no printed word, symbol, hypothesis or number of the counted statement or of its proof was altered. Layout and class: the article's own class is \texttt{article} with packages including geometry, graphicx, caption, subcaption, float, lmodern, fontenc, babel, mathtools, amsmath, amsthm, amssymb, amsfonts, mathrsfs; the wrapper is the lane's pinned \texttt{amsart} editorial wrapper at 10pt on A4, which restates the theorem-like environments the retained text needs and adds the article's own macro definitions that the retained text needs (\texttt{\textbackslash aut}, \texttt{\textbackslash calC}, \texttt{\textbackslash calL}, \texttt{\textbackslash d}, \texttt{\textbackslash hor}, \texttt{\textbackslash pd}, \texttt{\textbackslash ver}), with extra wrapper packages (bm, bbm, dsfont, xspace, cancel, mathtools, enumitem). The delivered numbering is the unit's own. Assets: no retained span contains a figure environment, a graphics-inclusion command, a PDF-page inclusion, a table or a TikZ picture; no image file is copied into this package, the package declares no asset dependency and no image file is redistributed with it. Licence: version arXiv:2402.16178v1 of this article is distributed under the Creative Commons Attribution 4.0 International licence, as recorded in the arXiv licence metadata for this exact version and shown on the abstract page https://arxiv.org/abs/2402.16178v1. A full rights notice is delivered at licenses/arxiv-2402.16178-CC-BY-4.0.txt. Characters: every retained excerpt is pure ASCII, so the delivered engine, XeLaTeX with its default Latin Modern text font, reports no missing character.
	\begin{proposition}\label{prop:basis_CASK}
		Let $(M,g,J,\nabla,\xi)$ be a CASK manifold. Then \begin{enumerate}[\normalfont(a)]
			\itemsep 0em
			\item The vector field $\xi$ is holomorphic and homothetic: $\mathcal L_\xi g = 2g$.
			\item The vector field $J\xi$ is holomorphic and Killing.
			\item The function $f=\frac{1}{2}g(\xi,\xi)$ is a K\"ahler potential for $g$ and a Hamiltonian function for $J\xi$, where our convention is $\d f=-\omega(J\xi,\cdot)$.
		\end{enumerate}
	\end{proposition}

		Since this is skew-symmetric, we see that $\d\iota_Zg_N=\pi^*\d\alpha=-2\pi^*\omega$. This shows that \begin{equation}\label{oH:eq}
			\omega_\hor=\begin{pmatrix}
				-\omega&0\\
				0&\omega^{-1}
			\end{pmatrix}.
		\end{equation}

	Note that \begin{equation}\label{pi*nabla:eq}
		(\pi^*\nabla)_Y\lambda=-\lambda\circ(\pi^*\nabla)(\pi^*X).
	\end{equation}

	\begin{proposition}\label{prop:moment_map}
		Let $X$ be an infinitesimal CASK automorphism and $Y$ its canonical lift as explained above. Then $$\mu_Y=\frac{1}{2}\left(g_N(Z,Y)+\pi^*\omega^{-1}(\lambda\circ(\pi^*\nabla)(\pi^*X),\lambda)\right)$$ is a $\omega_\hor$-Hamiltonian function for $Y$, where $Z=-\widetilde{J\xi}$. This assignment determines an equivariant (co)moment map $\mu:\aut(M)\longrightarrow\calC^\infty(N)$, $X\longmapsto\mu_X:=\mu_Y$, for the action of $\aut(M)$ on $N=T^*M$.
	\end{proposition}

	\begin{proof}
		We have to show that $\iota_Y\omega_\hor=-\d\mu_Y$. Recall that $\omega_\hor$ is given by \eqref{oH:eq} with respect to the splitting $T(T^*M)\cong\pi^*(TM)\oplus\pi^*(T^*M)$. Thus we have $$\iota_Y\omega_\hor=-\pi^*(\iota_X\omega)+\iota_{(\pi^*\nabla)_Y\lambda}\pi^*\omega^{-1}.$$
		
		Let us compute these two terms. First notice that, by Proposition~\ref{prop:basis_CASK}~(a) we have that 
		$\mathcal L_\xi\omega=2\omega$ and, hence $$2\iota_X\omega=\iota_X\mathcal L_\xi\omega =\iota_X\d(\iota_\xi\omega)=\mathcal L_X(\iota_\xi\omega)-\d(\iota_X\iota_\xi\omega)=\d(g(-J\xi,X)),$$ 
		since $\mathcal L_X(\iota_\xi\omega)=0$ (recall that $X$ is an infinitesimal CASK automorphism). So we see that \begin{equation}\label{iotaXomega:eq}
			\pi^*(\iota_X\omega)=\frac{1}{2}\d(g_N(Z,Y)).
		\end{equation}
	
		To compute the second term, let $A$ be an arbitrary section of $\pi^*(T^*M)\cong T^\ver(T^*M)$. Let $\{q^j\}$ be special real coordinates on $M$ and $\{q^j,p_j\}$ the corresponding canonical coordinates on $N$. With respect to these, we can write $A=A_i\d q^i$ and, using $(\pi^*\nabla)_Y\lambda=-\lambda\circ(\pi^*\nabla)(\pi^*X)=-p_i\pd{X^i}{q^j}\d q^j$, we get $$\pi^*\omega^{-1}((\pi^*\nabla)_Y\lambda,A)=-\omega^{jk}p_i\pd{X^i}{q^j}A_k.$$
		
		Thus, $\iota_{(\pi^*\nabla)_Y\lambda}\pi^*\omega^{-1}$ is the one-form on $T^*M$ which vanishes when applied to a horizontal vector field and evaluates on a vertical vector field $A=A_i\pd{}{p_i}$ (corresponding to the $A$ consider earlier by the canonical isomorphism $\pi^*(T^*M)\cong T^\ver(T^*M)$) as above. This means that we can write it, in local coordinates, as \begin{equation}\label{iotapi*nabla:eq}
			\iota_{(\pi^*\nabla)_Y\lambda}\pi^*\omega^{-1}=-\omega^{jk}p_i\pd{X^i}{q^j}\d p_k.
		\end{equation}
	
		Now, we compute the differential of our proposed moment map: $$\d\mu_Y=\frac{1}{2}\left(\d(g_N(Z,Y))+\d(\pi^*\omega^{-1}(\lambda\circ(\pi^*\nabla)(\pi^*X),\lambda))\right),$$ the second term of which can be computed (using \eqref{pi*nabla:eq} and \eqref{iotapi*nabla:eq}) in local coordinates as follows: \begin{align*}
			\frac{1}{2}\d(\pi^*\omega^{-1}(\lambda\circ(\pi^*\nabla)(\pi^*X),\lambda))&=\frac{1}{2}\d\left(\omega^{jk}p_ip_k\pd{X^i}{q^j}\right)=\frac{1}{2}\omega^{jk}\pd{X^i}{q^j}(p_i\d p_k+p_k\d p_i)\\
			&=\omega^{jk}p_i\pd{X^i}{q^j}\d p_k=-\iota_{(\pi^*\nabla)_Y\lambda}\pi^*\omega^{-1},
		\end{align*} where, in passing to the second line, we used that both $\omega^{jk}$ and $\pd{X^i}{q^j}$ are anti-symmetric in their indices; the latter fact is just the equation $\mathcal L_X \omega=0$ in ($\nabla$-affine) coordinates.\medskip
  
        This computation, together with \eqref{iotaXomega:eq}, concludes the proof that $\mu_Y$ is indeed a Hamiltonian function for $Y$.\medskip
		
		To show that the map $\mu$ is equivariant, let $X_1,X_2\in\aut(M)$ and $Y_1,Y_2\in\Gamma(TN)$ their corresponding canonical lifts. Note that the canonical lift of $[X_1,X_2]$ is precisely $[Y_1,Y_2]$. Using this and the fact that the moment map is constructed in a canonical way, we obtain $$\calL_{Y_1}(\mu_{Y_2})=\frac{1}{2}\left(g_N(Z,[Y_1,Y_2])+\pi^*\omega^{-1}(\lambda\circ(\pi^*\nabla)(\pi^*([X_1,X_2])),\lambda)\right)=\mu_{[Y_1,Y_2]}.$$
	\end{proof}

\end{document}
